\section{Introduction}\label{sec:intro}

Diabetic retinopathy damages the small blood vessels of the retina and is a leading preventable cause of
vision loss. It is diagnosed from colour fundus photographs, and the first computational step in most automated
screening systems is segmenting the retinal vessels~\cite{fraz2012}. India alone had an estimated 89.8 million
adults with diabetes in 2024~\cite{idf2025}. Screening camps capture fundus images in bulk, often on modest
hardware and with few specialists available to read them, so throughput (images per second) is a practical
requirement.

The multi-scale Frangi vesselness filter~\cite{frangi1998} is a widely used classical method for this step. It
needs no training data, behaves deterministically and can be written in plain C. Its cost grows with image
resolution and with the number of scales, because every pixel needs several Gaussian-derivative convolutions
followed by an eigenvalue computation. A complete screening pipeline adds contrast enhancement before the filter
and binary clean-up after it, and these stages stress a processor in different ways: convolution is
memory-bound, the eigenvalue step is arithmetic-heavy, CLAHE builds histograms, and morphology moves single
bytes.

Multicore CPUs are the hardware every screening site already has. A developer choosing how to parallelise such
a pipeline needs to know, stage by stage, where OpenMP threads help, where they stop helping, and why.
Published parallel retinal work has used other algorithms or isolated single stages~\cite{arguello2018,bensayadia2020}, and the closest cross-platform study covers a single enhancement filter~\cite{song2022}.

This paper implements the full pipeline (green-channel extraction, CLAHE, multi-scale Frangi filtering,
thresholding and morphological clean-up) in C, parallelises every stage with OpenMP, and analyses its
performance stage by stage against a pragma-free serial build of the same code. Accuracy is measured on the
DRIVE benchmark~\cite{staal2004}, with STARE~\cite{hoover2000} as a cross-dataset check, using parameters tuned
on the DRIVE training images only. The contributions are:
\begin{itemize}
\item a documented serial/OpenMP C implementation whose outputs are bit-identical across all thread counts and
design variants, verified on every test image;
\item a stage-by-stage measurement of time, speedup, efficiency and Karp--Flatt serial fraction on a
2-core/4-thread cloud CPU, including the effect of Hyper-Threading and thread placement;
\item an ablation of three alternative OpenMP designs (dynamic tile scheduling, atomic histograms and
scale-level parallelism) and a resolution-scaling study from 0.33 to 21.1 megapixels;
\item segmentation accuracy on DRIVE and STARE reported alongside published lightweight deep models.
\end{itemize}

The study tests four hypotheses. H1 and H4 are taken unchanged from the project proposal. The proposal's two
GPU hypotheses fall outside this OpenMP-only study and are replaced by H2 and H3, which were written after a
short pilot run (two images, three repetitions) used to debug the harness; the pilot data are not reported.
\textbf{H1}: OpenMP speedup grows sub-linearly with thread count and is limited by the number of physical
cores and by memory bandwidth.
\textbf{H2}: OpenMP speedup is stable or grows with image size, because the fixed costs of creating parallel
regions shrink relative to the work.
\textbf{H3}: the best decomposition differs by stage: owner-computes tile histograms beat atomic
histograms in CLAHE, and parallelism inside each scale beats running the scales concurrently in the Frangi
stage, because the four scales have unequal cost.
\textbf{H4}: the classical pipeline does not match the best lightweight deep models in Dice (for example,
0.8070 on DRIVE for LightVesselNet~\cite{sobhan2026}); its value is that it is training-free and its parallel
behaviour can be explained.

Section~\ref{sec:related} reviews related work, Section~\ref{sec:method} describes the pipeline, the OpenMP
design and the measurement protocol, Section~\ref{sec:results} reports the results and
Section~\ref{sec:discussion} discusses them against the hypotheses.
